% GENERATED FILE. Edit canonical lessons, then run npm run generate:books.
% canonical-source-hash: fnv1a64:57fc7093094d6478
% canonical-lessons: IT-C412-vigile, IT-C412-pedone, IT-C412-strisce, IT-C412-investire, IT-C412-contravvenzione, IT-R412-ripasso-405-412

\chapter{Pedoni e vigili}
\label{ch:it-pedoni-e-vigili}
\hlchaptermodality{412}

\begin{chapteropening}
\textbf{By the end of this chapter:} \emph{I can talk about road rules in town: the traffic warden, pedestrians, the zebra crossing, running someone over and fines.}

Complete the last lesson of chapter 412: I can talk about road rules in town: the traffic warden, pedestrians, the zebra crossing, running someone over and fines.
\end{chapteropening}

\section[il vigile]{il vigile — the traffic warden, the local police officer}

\label{lesson:IT-C412-vigile}



\begin{quote}
Before the new one: say the Italian for the light, then the Italian for to unscrew.
\end{quote}

\subsection*{You'll want to know: il vigile}
\textbf{il vigile} — \textquotedblleft{}the traffic warden, the local police officer\textquotedblright{}.

\begin{quote}
Il vigile ci ha fermato all'incrocio.
\end{quote}

\begin{quote}
\emph{The warden stopped us at the crossroads.}
\end{quote}

In full \emph{il vigile urbano}; \emph{i vigili del fuoco} are firefighters.

\begin{cousinweb}
\emph{vigile} is a learned borrowing from Latin \emph{vigil} 'awake, watchful' (Lewis \& Short): the one who keeps watch.
\end{cousinweb}

\begin{grammarlens}[title={What you've built}]
One more word for telling, explaining and describing at B1.
\end{grammarlens}

\subsection*{Your turn}
\begin{itemize}
\raggedright
  \item \emph{Say it:} \emph{il vigile}
  \item \emph{Say it:} \emph{il vigile}, once more
  \item \emph{Say it:} a sentence of your own with \emph{il vigile}
  \item \emph{Recall:} say the Italian for the light, then the Italian for to unscrew, then say \emph{il vigile} again
\end{itemize}

\begin{tcolorbox}[breakable,colback=teal!4,colframe=teal!35!black,title={Before you move on}]
What does il vigile mean? (\textquotedblleft{}The traffic warden, the local police officer\textquotedblright{}.) Say it once more.
\end{tcolorbox}
\section[il pedone]{il pedone — the pedestrian}

\label{lesson:IT-C412-pedone}



\begin{quote}
Before the new one: say the Italian for to unscrew, then the Italian for the traffic warden.
\end{quote}

\subsection*{You'll want to know: il pedone}
\textbf{il pedone} — \textquotedblleft{}the pedestrian\textquotedblright{}.

\begin{quote}
Attenzione ai pedoni!
\end{quote}

\begin{quote}
\emph{Watch out for pedestrians!}
\end{quote}

\emph{La zona pedonale} is a pedestrian zone. In chess, \emph{il pedone} is a pawn.

\begin{cousinweb}
\emph{pedone} is inherited from Late Latin \emph{pedō} 'foot soldier', built on \emph{pēs, pedis} 'foot' (EtymDB).
\end{cousinweb}

\begin{grammarlens}[title={What you've built}]
One more word for telling, explaining and describing at B1.
\end{grammarlens}

\subsection*{Your turn}
\begin{itemize}
\raggedright
  \item \emph{Say it:} \emph{il pedone}
  \item \emph{Say it:} \emph{il pedone}, once more
  \item \emph{Say it:} a sentence of your own with \emph{il pedone}
  \item \emph{Recall:} say the Italian for to unscrew, then the Italian for the traffic warden, then say \emph{il pedone} again
\end{itemize}

\begin{tcolorbox}[breakable,colback=teal!4,colframe=teal!35!black,title={Before you move on}]
What does il pedone mean? (\textquotedblleft{}The pedestrian\textquotedblright{}.) Say it once more.
\end{tcolorbox}
\section[le strisce]{le strisce — the zebra crossing; the stripes}

\label{lesson:IT-C412-strisce}



\begin{quote}
Before the new one: say the Italian for the traffic warden, then the Italian for the pedestrian.
\end{quote}

\subsection*{You'll want to know: le strisce}
\textbf{le strisce} — \textquotedblleft{}the zebra crossing; the stripes\textquotedblright{}.

\begin{quote}
Attraversa sulle strisce, è più sicuro.
\end{quote}

\begin{quote}
\emph{Cross at the zebra crossing, it is safer.}
\end{quote}

Singular \emph{la striscia}, a stripe; in full \emph{le strisce pedonali}.

\begin{cousinweb}
\emph{strisce} is the plural of \emph{striscia} 'stripe', which EtymDB traces to Latin \emph{stria} 'furrow, groove'; it is inherited.
\end{cousinweb}

\begin{grammarlens}[title={What you've built}]
One more word for telling, explaining and describing at B1.
\end{grammarlens}

\subsection*{Your turn}
\begin{itemize}
\raggedright
  \item \emph{Say it:} \emph{le strisce}
  \item \emph{Say it:} \emph{le strisce}, once more
  \item \emph{Say it:} a sentence of your own with \emph{le strisce}
  \item \emph{Recall:} say the Italian for the traffic warden, then the Italian for the pedestrian, then say \emph{le strisce} again
\end{itemize}

\begin{tcolorbox}[breakable,colback=teal!4,colframe=teal!35!black,title={Before you move on}]
What does le strisce mean? (\textquotedblleft{}The zebra crossing; the stripes\textquotedblright{}.) Say it once more.
\end{tcolorbox}
\section[investire]{investire — to run over; to invest}

\label{lesson:IT-C412-investire}



\begin{quote}
Before the new one: say the Italian for the pedestrian, then the Italian for the zebra crossing.
\end{quote}

\subsection*{You'll want to know: investire}
\textbf{investire} — \textquotedblleft{}to run over; to invest\textquotedblright{}.

\begin{quote}
Una macchina ha investito un gatto.
\end{quote}

\begin{quote}
\emph{A car ran over a cat.}
\end{quote}

With money it means to invest: \emph{investire in una casa}.

\begin{cousinweb}
\emph{investire} is a learned borrowing from Latin \emph{investīre} 'to clothe, to cover' (EtymDB), built on \emph{vestis} 'clothing'.
\end{cousinweb}

\begin{grammarlens}[title={What you've built}]
One more word for telling, explaining and describing at B1.
\end{grammarlens}

\subsection*{Your turn}
\begin{itemize}
\raggedright
  \item \emph{Say it:} \emph{investire}
  \item \emph{Say it:} \emph{investire}, once more
  \item \emph{Say it:} a sentence of your own with \emph{investire}
  \item \emph{Recall:} say the Italian for the pedestrian, then the Italian for the zebra crossing, then say \emph{investire} again
\end{itemize}

\begin{tcolorbox}[breakable,colback=teal!4,colframe=teal!35!black,title={Before you move on}]
What does investire mean? (\textquotedblleft{}To run over; to invest\textquotedblright{}.) Say it once more.
\end{tcolorbox}
\section[la contravvenzione]{la contravvenzione — the fine, the traffic ticket}

\label{lesson:IT-C412-contravvenzione}



\begin{quote}
Before the new one: say the Italian for the zebra crossing, then the Italian for to run over.
\end{quote}

\subsection*{You'll want to know: la contravvenzione}
\textbf{la contravvenzione} — \textquotedblleft{}the fine, the traffic ticket\textquotedblright{}.

\begin{quote}
Il vigile mi ha fatto una contravvenzione.
\end{quote}

\begin{quote}
\emph{The warden gave me a ticket.}
\end{quote}

Formal; in speech people say \emph{la multa}.

\begin{cousinweb}
\emph{contravvenzione} is a learned borrowing from Medieval Latin \emph{contrāventiō}, built on \emph{contrā} 'against' and \emph{venīre} 'to come': coming against a rule.
\end{cousinweb}

\begin{grammarlens}[title={What you've built}]
One more word for telling, explaining and describing at B1.
\end{grammarlens}

\subsection*{Your turn}
\begin{itemize}
\raggedright
  \item \emph{Say it:} \emph{la contravvenzione}
  \item \emph{Say it:} \emph{la contravvenzione}, once more
  \item \emph{Say it:} a sentence of your own with \emph{la contravvenzione}
  \item \emph{Recall:} say the Italian for the zebra crossing, then the Italian for to run over, then say \emph{la contravvenzione} again
\end{itemize}

\begin{tcolorbox}[breakable,colback=teal!4,colframe=teal!35!black,title={Before you move on}]
What does la contravvenzione mean? (\textquotedblleft{}The fine, the traffic ticket\textquotedblright{}.) Say it once more.
\end{tcolorbox}
\section[(dialogue)]{Strade, stazioni e alberghi}

\label{lesson:IT-R412-ripasso-405-412}



\begin{quote}
Say the Italian for the toll booth and the lane.
\end{quote}

\subsection*{Your turn}
Say these aloud:

\begin{itemize}
\raggedright
  \item the toll booth, the lane, the ring road, the junction, to overtake
  \item to speed up, the brake, the braking, to steer, to skid
  \item the windscreen, the windscreen wiper, the wing mirror, the horn, to mist up
  \item the petrol station attendant, the diesel, the fuel, the tank, to refuel
  \item to get a puncture, the puncture, the tyre, the jack, to pump up
  \item the breakdown, to break down, the tow truck, to tow, the garage
  \item the bonnet, the boot, the clutch, the light, to unscrew
  \item the traffic warden, the pedestrian, the zebra crossing, to run over, the fine
\end{itemize}

\begin{tcolorbox}[breakable,colback=teal!4,colframe=teal!35!black,title={Before you move on}]
The toll booth? (\textbf{il casello}.) The lane? (\textbf{la corsia}.) The ring road? (\textbf{la tangenziale}.) The junction? (\textbf{lo svincolo}.) To overtake? (\textbf{sorpassare}.) To speed up? (\textbf{accelerare}.) The brake? (\textbf{il freno}.) The braking? (\textbf{la frenata}.) To steer? (\textbf{sterzare}.) To skid? (\textbf{sbandare}.) The windscreen? (\textbf{il parabrezza}.) The windscreen wiper? (\textbf{il tergicristallo}.) The wing mirror? (\textbf{lo specchietto}.) The horn? (\textbf{il clacson}.) To mist up? (\textbf{appannarsi}.) The petrol station attendant? (\textbf{il benzinaio}.) The diesel? (\textbf{il gasolio}.) The fuel? (\textbf{il carburante}.) The tank? (\textbf{il serbatoio}.) To refuel? (\textbf{rifornire}.) To get a puncture? (\textbf{forare}.) The puncture? (\textbf{la foratura}.) The tyre? (\textbf{lo pneumatico}.) The jack? (\textbf{il cric}.) To pump up? (\textbf{gonfiare}.) The breakdown? (\textbf{il guasto}.) To break down? (\textbf{guastarsi}.) The tow truck? (\textbf{il carro attrezzi}.) To tow? (\textbf{rimorchiare}.) The garage? (\textbf{l'officina}.) The bonnet? (\textbf{il cofano}.) The boot? (\textbf{il bagagliaio}.) The clutch? (\textbf{la frizione}.) The light? (\textbf{il fanale}.) To unscrew? (\textbf{svitare}.) The traffic warden? (\textbf{il vigile}.) The pedestrian? (\textbf{il pedone}.) The zebra crossing? (\textbf{le strisce}.) To run over? (\textbf{investire}.) The fine? (\textbf{la contravvenzione}.)
\end{tcolorbox}
